\section{Introduction}\label{sec:intro}

Income distributions are usually drawn as a single hump with a long right tail, and most summary measures of inequality and polarisation are built on that picture \citep{Wolfson1994,EstebanRay1994,DuclosEstebanRay2004}. There are good reasons to expect exceptions among people whose income is set largely by rule. A statutory benefit paid at a flat rate places everyone who receives it and has little else at the same point, and a population in which a large share lives on such a benefit can have a mode at the benefit amount that is distinct from the mode formed by everyone else. Whether this happens, and where, is an empirical question. The literature on minimum income protection measures the level of statutory floors against poverty lines \citep{Bahle2011,Immervoll2012,MarxNelson2013,Goedeme2013,MarchalMarxVanMechelen2014}, and the literature on bunching measures the response of taxable income to kinks and notches \citep{Saez2010,ChettyFriedmanOlsenPistaferri2011,KlevenWaseem2013,Kleven2016}, but neither asks what floors do to the shape of the income distribution as a whole.

This paper studies that question for people who live alone, first in Denmark and Norway and then in eight other countries, using the harmonised microdata of the Luxembourg Income Study \citep{LIS2024,GornickJantti2013}. One-person households are the natural place to look. Their income is not pooled with anyone else's, so a benefit paid to an individual is the household's income; they are disproportionately old and disproportionately young, the two groups whose incomes are most directly set by statute; and in Denmark and Norway they are numerous enough, and the register-based LIS files large enough, that differences in the shape of the distribution can be tested rather than eyeballed. Shape is tested with a test of unimodality for binned incomes that is valid on the boundary of the null hypothesis, developed in a related paper \citep{Orujov2026}; where it rejects, it also locates the valley and the two modes, so that their positions can be compared with the statutory amounts.

The main finding concerns Danes aged 65 and over who live alone. In each of the three waves from 2013 onwards their disposable income is bimodal, with a lower mode at between 125{,}000 and 140{,}000 kroner and an upper mode about 35{,}000 kroner higher. Among those who live almost entirely on public transfers, roughly half of all older people living alone, the distribution is bimodal in every wave from 2004 onwards, including the earlier waves in which the sample is a quarter of the size. The two modes are not the floor and the rest. They are two floors. The public old-age pension, the \emph{folkepension}, paid at the full rate with its supplement to a single person, produces one; the same pension plus the pensioners' housing benefit, \emph{boligydelse}, produces the other. Within each of the two groups income is unimodal in every wave, once former civil servants and others with sizeable public contributory pensions are set aside; the distance between the medians of the two groups is close to the median benefit; one recipient in six or seven receives exactly the statutory maximum benefit in the two years in which this can be checked; and removing the benefit from income removes the second mode. The recipients are, almost without exception, renters, and two thirds of the non-recipients own their home. Among older people with occupational or private pensions, and among people aged 30 to 64, there is little or no evidence against unimodality once a small mass of near-zero incomes is set aside.

Norway supplies a contrast. It too guarantees a minimum pension to a single person, and because minimum pensioners pay essentially no income tax, the disposable incomes of Norwegians who live on it sit close to the statutory amount. But only four to seven per cent of Norwegians aged 67 and over who live alone depend on it, because the earnings-related national pension lifts almost everyone above the floor, and the distribution of their disposable income is unimodal in every register wave from 2007 to 2022. A floor imprints a mode only when a large share of a group sits on it.

The Danish and Norwegian cases point to three conditions: a large share of the group lives on a statutory floor; a second transfer is paid to a substantial part of that group but not to all of it; and the transfer is large relative to the spread of income within each part. Section~\ref{sec:countries} tests them in eight other countries, pooling all LIS waves from 2010 to 2021 to obtain samples large enough for the test. The structure reappears where the conditions hold. In the Netherlands the flat-rate state pension and the housing allowance produce modes at the floor and 24 per cent above it; in the United Kingdom the state pension and Housing Benefit produce two modes; and in France the old-age minimum and the housing allowance did so from 2010 to 2014, while recipients were about 40 per cent richer than non-recipients, and ceased to as that gap closed. In each case subtracting the allowance removes the second mode. In Australia, where only a sixth of pensioners living alone receive Rent Assistance, the evidence is weak, as a Danish subsampling experiment says it must be. In Sweden, Germany, the United States and Ireland, where a condition fails, no second mode is due to the allowance.

The Danish student grant, SU, paid at a flat monthly rate to students in higher education, leaves a mark of a different kind at the other end of the life course: the gross income of Danes under 30 who live alone bunches at twelve times the monthly rate in every wave, and in the two largest waves their disposable income has a mode just below the grant (Appendix~\ref{app:su}).

The feature documented here is, as far as I know, new, and it has a direct consequence for measurement: the conventional at-risk-of-poverty line, 60 per cent of median equivalised income, falls in the valley between the two modes in the recent waves, so that nine in ten transfer-reliant older people without the housing benefit are below it and nine in ten of those with it are above. The method, a test of unimodality that locates the modes, read against the statutory amounts and combined with a decomposition by income source, can be applied wherever flat-rate transfers make up a large part of the income of some group.

Section~\ref{sec:inst} describes the Danish and Norwegian pension floors, the Danish housing benefit and the student grant. Section~\ref{sec:data} describes the data and the two statistical tools. Section~\ref{sec:results} presents the results for Denmark, Norway and eight other countries, and Section~\ref{sec:disc} discusses what they imply for the measurement of low income. The appendix reports Sweden and Finland, whose LIS files are samples in which the test has limited power against differences of the size found in Denmark, and the student grant.
